% Source: https://artofproblemsolving.com/community/c3h2157281_counting_and_probability_questions
% Stable identity: aops-topic-2157281-q07
If $n=p_1^{e_1}\cdots p_k^{e_k}$, then $n$ has $(e_1+1)\cdots(e_k+1)$ divisors. Writing $12$ as such a product:
\begin{itemize}
\item $12$: $n=p^{11}\ge2^{11}$, too large;
\item $6\cdot2$: $n=p^5q$: $2^5\cdot3=96$ ($2^5\cdot5=160$ and $3^5\cdot2=486$ are too large);
\item $4\cdot3$: $n=p^3q^2$: $2^3\cdot3^2=72$ ($2^3\cdot5^2=200$, $3^3\cdot2^2=108$ are too large);
\item $3\cdot2\cdot2$: $n=p^2qr$: $2^2\cdot3\cdot5=60$, $2^2\cdot3\cdot7=84$, $3^2\cdot2\cdot5=90$ (the next ones, $2^2\cdot3\cdot11=132$, $2^2\cdot5\cdot7=140$, $3^2\cdot2\cdot7=126$, $5^2\cdot2\cdot3=150$, are too large).
\end{itemize}
The numbers are $60,72,84,90,96$, with sum $\boxed{402}$.
